% =====================================================================
%  10. Seuils normatifs et validation des soudures
% =====================================================================
\clearpage
\section{Seuils normatifs et validation des soudures}
\label{sec:validation}

Les sections précédentes prédisent quatre propriétés. Pour juger la
\og qualité \fg{} d'une soudure, il faut encore décider si ces propriétés sont
\emph{suffisantes}. Cette section définit des seuils par cible à partir de
normes publiées, puis une règle de décision qui tient compte de l'erreur des
modèles. Le profil et le niveau de prudence ont été choisis par l'équipe ;
la recherche complète et ses sources sont dans
\fichier{reports/18\_seuils\_validation\_fr.md}, le code dans
\notebook{09\_validation.ipynb}.

\subsection{Référence normative}

La base décrit surtout du \textbf{métal fondu} : 1\,121 des 1\,555 lignes de la
cohorte proviennent du soudage manuel à l'électrode enrobée (MMA). Les normes de
\emph{classification des produits consommables} testent précisément ce métal
fondu pur ; ce sont donc les références directes. Les normes de qualification de
modes opératoires (ISO~15614-1, DNV-OS-C401) testent un joint et renvoient au
métal de base : elles fixent le niveau d'exigence d'une application, pas une
valeur unique.

\paragraph{ISO 2560-A~\cite{iso_2560,twi_consumables}.} Pour les aciers non
alliés et à grains fins, une classe impose simultanément une limite d'élasticité
minimale, une plage de résistance et un allongement minimal
(\cref{tab:iso2560}). La résilience est notée par la température à laquelle la
\textbf{moyenne de trois éprouvettes Charpy atteint 47\,J} (symbole A : $+20$\,°C,
0 : 0\,°C, 2 : $-20$\,°C, 3 : $-30$\,°C, 4 : $-40$\,°C, 5 : $-50$\,°C,
6 : $-60$\,°C) ; une seule valeur individuelle peut être inférieure à 47\,J, sans
descendre sous 32\,J. La limite d'élasticité est $R_{eL}$ (ou $R_{p0,2}$ sans
palier) et l'allongement est mesuré sur une longueur de cinq diamètres.

\begin{table}[htbp]
\centering
\small
\caption{ISO 2560-A, tableau 1A : propriétés de traction du métal fondu par
classe, et ISO 18275-A pour les aciers à haute résistance.}
\label{tab:iso2560}
\begin{tabular}{@{}l c r c r@{}}
\toprule
Norme & Symbole & \YS{} min (MPa) & \UTS{} (MPa) & Allongement min (\%) \\
\midrule
ISO 2560-A & 35 & 355 & 440--570 & 22 \\
           & 38 & 380 & 470--600 & 20 \\
           & \textbf{42} & \textbf{420} & \textbf{500--640} & \textbf{20} \\
           & 46 & 460 & 530--680 & 20 \\
           & 50 & 500 & 560--720 & 18 \\
\midrule
ISO 18275-A & 55 & 550 & 610--780 & 18 \\
            & 62 & 620 & 690--890 & 18 \\
            & 69 & 690 & 760--960 & 17 \\
\bottomrule
\end{tabular}
\end{table}

\paragraph{Autres références.} ISO 18275-A~\cite{iso_18275} prolonge le même
système (47\,J) au-delà de 500\,MPa. L'AWS A5.1 utilise un système à 27\,J :
une électrode E7018 doit atteindre \YS{} $\geq$ 400\,MPa, \UTS{} $\geq$ 490\,MPa,
22\,\% d'allongement et 27\,J à $-29$\,°C~\cite{aws_a51}. En offshore, la
DNV exige pour le métal fondu la même énergie que pour le métal de base en sens
travers, à la température de la nuance, avec une valeur individuelle au moins
égale à 70\,\% de la moyenne requise~\cite{dnv_os_c401} ; pour le métal de
base, l'énergie en sens long vaut environ \YS/10 en joules (42\,J pour 420\,MPa,
69\,J pour 690\,MPa)~\cite{dnv_os_b101}.

\paragraph{Compromis.} Les seuils ne sont pas indépendants.
\begin{itemize}
  \item \textbf{Résistance contre allongement} : les normes abaissent l'allongement minimal quand la classe monte (22 à 15\,\%) ; dans nos données, la corrélation de Spearman entre \YS{} et allongement vaut $-0{,}73$ (622 soudures, 33 séries).
  \item \textbf{Résistance contre ténacité} : un métal fondu se rend facilement résistant, plus difficilement tenace ; la composition (Mn faible, Ni élevé) peut déplacer ce compromis sans le supprimer~\cite{keehan_2001}. Sur des monopieux d'éoliennes en S355, l'énergie Charpy du métal fondu est environ 30\,\% plus faible que celle du métal de base~\cite{mehmanparast_2018}.
  \item \textbf{Sur-résistance contre borne haute} : la plage d'\UTS{} limite la sur-résistance ; 27\,\% des \UTS{} de la cohorte dépassent 640\,MPa.
\end{itemize}

\subsection{Profil retenu et règle de décision}

\begin{decision}
Le profil retenu est \textbf{ISO 2560-A E~42~2} : \YS{} $\geq$ 420\,MPa,
\UTS{} $\geq$ 500\,MPa, allongement $\geq$ 20\,\% et énergie Charpy moyenne
$\geq$ 47\,J à $-20$\,°C, tous \textbf{critiques} ; \UTS{} $\leq$ 640\,MPa est
\textbf{informatif}. La classe 42 est au centre des données (médianes : \YS{}
496\,MPa, \UTS{} 576\,MPa, allongement 26,6\,\%) et $-20$\,°C est l'une des
températures d'essai les plus fréquentes (76 couples). D'autres profils sont
documentés (AWS E7018, E~42~4 à $-40$\,°C pour l'offshore et l'éolien,
ISO 18275 E~55~4) mais non retenus.
\end{decision}

Une prédiction a une erreur. Pour en tenir compte, nous appliquons la règle de
la \emph{bande de garde} du JCGM~106 (§8.3)~\cite{jcgm_106}. Pour un critère de
seuil bas $L$ (resp. haut $H$), une valeur $\hat{y}$ d'incertitude type $u$
reçoit le statut :
\begin{equation}
  \text{C (conforme) si } \hat{y} - w \geq L, \qquad
  \text{NC (non conforme) si } \hat{y} + w < L, \qquad
  \text{I (indéterminé) sinon},
\end{equation}
avec $w = k\,u$ et symétriquement pour $H$. Avec $k = 2$ ($w = U = 2u$, règle
par défaut de l'ISO~14253-1), la probabilité d'accepter une soudure non conforme
est au plus de 2,3\,\% pour une erreur normale. La décision globale est alors :
\begin{itemize}
  \item \textbf{VALIDÉ} si tous les critères sont C ;
  \item \textbf{NON VALIDÉ} si au moins un critère critique est NC ;
  \item \textbf{PARTIEL} sinon : la soudure doit être confirmée par un essai réel.
\end{itemize}

\subsection{Mise en œuvre}

Le notebook 09 applique cette règle à deux sources sur les 1\,555 soudures de la
cohorte :
\begin{itemize}
  \item les \textbf{mesures} de la base, avec $w = 0$ (aucune incertitude de mesure n'est publiée) ; une cible absente donne I, et le Charpy n'est tranché que si l'unique couple $E(T)$ suffit, l'énergie étant supposée croissante avec la température ;
  \item les \textbf{prédictions hors pli} de la sélection à six familles, pour chaque protocole : chaque soudure est prédite par le modèle externe qui ne l'a pas vue (ni sa série en groupé), et l'énergie Charpy est prédite à $-20$\,°C.
\end{itemize}
L'incertitude $u$ est la RMSE hors pli réunie de chaque protocole et cible. Le
\cref{tab:limites} donne les limites d'acceptation qui en résultent. Les
40 modèles sont vérifiés par empreinte, reproduisent exactement leurs
prédictions exportées, et six tests couvrent la règle.

\begin{table}[htbp]
\centering
\small
\caption{Bande de garde $w = 2u$ et limites appliquées à la valeur prédite
(C si / NC si), profil E~42~2.}
\label{tab:limites}
\setlength{\tabcolsep}{4pt}
\begin{tabular}{@{}l r l r l@{}}
\toprule
 & \multicolumn{2}{c}{Validation aléatoire} & \multicolumn{2}{c}{Validation groupée} \\
\cmidrule(lr){2-3}\cmidrule(lr){4-5}
Critère & $w$ & C si / NC si & $w$ & C si / NC si \\
\midrule
\YS{} $\geq$ 420\,MPa          & 79,1 & $\geq 499{,}1$ / $< 340{,}9$ & 144,5 & $\geq 564{,}5$ / $< 275{,}5$ \\
\UTS{} $\geq$ 500\,MPa         & 73,2 & $\geq 573{,}2$ / $< 426{,}8$ & 141,4 & $\geq 641{,}5$ / $< 358{,}6$ \\
\UTS{} $\leq$ 640\,MPa         & 73,2 & $\leq 566{,}8$ / $> 713{,}2$ & 141,4 & $\leq 498{,}6$ / $> 781{,}5$ \\
Allongement $\geq$ 20\,\%      & 5,8  & $\geq 25{,}8$ / $< 14{,}2$   & 8,3   & $\geq 28{,}3$ / $< 11{,}7$ \\
KV($-20$\,°C) $\geq$ 47\,J     & 39,5 & $\geq 86{,}5$ / $< 7{,}5$    & 96,0  & $\geq 143{,}0$ / jamais NC \\
\bottomrule
\end{tabular}
\end{table}

La plage d'\UTS{} (140\,MPa) est plus étroite que deux bandes de garde dans les
deux protocoles : aucune prédiction ne peut rendre conformes à la fois
\UTS{} $\geq$ 500 et \UTS{} $\leq$ 640\,MPa.

\subsection{Résultats}

\begin{figure}[htbp]
\centering
\includegraphics[width=\textwidth]{val_decisions}
\caption{Décision de validation des 1\,555 soudures, profil ISO 2560-A E~42~2,
$k = 2$ : mesures de la base et prédictions hors pli des deux protocoles.}
\label{fig:val_decisions}
\end{figure}

\begin{table}[htbp]
\centering
\small
\caption{Décisions sur la cohorte (règle retenue).}
\label{tab:val_decisions}
\begin{tabular}{@{}l rrr@{}}
\toprule
Source & VALIDÉ & PARTIEL & NON VALIDÉ \\
\midrule
Mesures                 & 51 (3,3\,\%) & 1\,274 (81,9\,\%) & 230 (14,8\,\%) \\
Modèle KFold aléatoire  & 0            & 1\,538 (98,9\,\%) & 17 (1,1\,\%) \\
Modèle GroupKFold       & 0            & 1\,555 (100\,\%)  & 0 \\
\bottomrule
\end{tabular}
\end{table}

\begin{table}[htbp]
\centering
\small
\caption{Accord entre statut prédit et statut mesuré, sur les lignes où la
mesure tranche. Part tranchée : prédictions C ou NC ; fausses acceptations :
prédit C, mesuré NC, rapporté aux C prédits.}
\label{tab:val_accord}
\begin{tabular}{@{}l r ll ll@{}}
\toprule
 & & \multicolumn{2}{c}{Validation aléatoire} & \multicolumn{2}{c}{Validation groupée} \\
\cmidrule(lr){3-4}\cmidrule(lr){5-6}
Critère & Lignes & Part tranchée & Fausses acc. & Part tranchée & Fausses acc. \\
\midrule
\YS{} $\geq$ 420     & 706 & 51\,\% & 2/353 (0,6\,\%) & 16\,\% & 2/112 (1,8\,\%) \\
\UTS{} $\geq$ 500    & 677 & 52\,\% & 0/350           & 19\,\% & 0/130 \\
\UTS{} $\leq$ 640    & 677 & 56\,\% & 4/307 (1,3\,\%) & 9\,\%  & 1/43 (2,3\,\%) \\
Allongement          & 631 & 60\,\% & 5/376 (1,3\,\%) & 27\,\% & 1/173 (0,6\,\%) \\
KV($-20$\,°C)        & 465 & 86\,\% & 1/399 (0,3\,\%) & 7\,\%  & 0/34 \\
\bottomrule
\end{tabular}
\end{table}

\paragraph{Lecture.}
\begin{itemize}
  \item \textbf{Mesures.} Seules 141 lignes ont les quatre cibles mesurées : 1\,253 des 1\,274 PARTIEL ont au moins une cible absente. Les 230 NON VALIDÉ échouent sur \YS{} (109), \UTS{} minimale (83), l'allongement (72) ou le Charpy (44), avec des recouvrements.
  \item \textbf{Modèles.} Aucun VALIDÉ, conséquence directe de la plage d'\UTS{} vide (\cref{tab:limites}). Les 17 NON VALIDÉ du modèle aléatoire ne sont jamais contredits par la mesure : 7 sont confirmés, 10 ont une mesure incomplète.
  \item \textbf{Calibration.} Les fausses acceptations par critère restent sous la borne de 2,3\,\% dans les deux protocoles, et un rejet prédit sur une borne basse n'est jamais contredit par la mesure : la bande $k = 2$ est prudente sur ces données.
  \item \textbf{Pouvoir de décision.} Le modèle aléatoire tranche la moitié des critères de traction et 86\,\% des critères Charpy ; le modèle groupé n'en tranche que 7 à 27\,\%.
  \item \textbf{Sensibilité.} Si l'on ignore \UTS{} $\leq$ 640\,MPa, le modèle aléatoire valide 125 soudures : la mesure en confirme 5, en contredit 2 (1,6\,\%) et ne permet pas de conclure pour 118. Le modèle groupé n'en valide toujours aucune. Rendre ce critère critique ajoute surtout des NON VALIDÉ.
\end{itemize}

\begin{precaution}
Cette évaluation est rétrospective, sur les données qui ont servi à choisir les
modèles. La RMSE est une erreur globale par cible ; pour Charpy, elle est
estimée aux températures observées et appliquée à $-20$\,°C. Les règles sur les
valeurs individuelles des éprouvettes ne peuvent pas être vérifiées avec une
énergie moyenne prédite. Enfin, ces seuils concernent le métal fondu, pas un
joint soudé complet. Des intervalles de prédiction propres à chaque soudure
(par exemple CV+~\cite{barber_2021}) affineraient la règle.
\end{precaution}

\begin{aretenir}
Avec le profil ISO 2560-A E~42~2 et une bande de garde de deux erreurs types,
le modèle \textbf{ne valide aucune soudure} : la plage d'\UTS{} est plus étroite
que l'incertitude. Ses décisions tranchées sont fiables (moins de 2,3\,\% de
fausses acceptations par critère), mais rares pour une série nouvelle. Le modèle
sert donc à \textbf{trier des candidats avant essai}, pas à qualifier une
soudure ; le traitement de la borne haute d'\UTS{} reste un choix de risque.
\end{aretenir}
